\documentclass[12pt]{article}
\usepackage[utf8]{inputenc}
\usepackage{amsmath, amssymb, amsthm}
\usepackage{geometry}

\geometry{a4paper, margin=1in}

\newtheorem{theorem}{Theorem}
\newtheorem{lemma}[theorem]{Lemma}
\newtheorem{proposition}[theorem]{Proposition}
\theoremstyle{definition}
\newtheorem{definition}[theorem]{Definition}
\theoremstyle{remark}
\newtheorem{remark}[theorem]{Remark}

\title{Project Euler Problem 12: Highly Divisible Triangular Number}
\author{}
\date{}

\begin{document}
\maketitle

\section{Problem Statement}

Find the smallest triangular number $T_n = \frac{n(n+1)}{2}$ with more than 500 positive divisors.

\section{Formal Development}

\begin{definition}
For $m \in \mathbb{Z}^+$ with canonical factorization $m = \prod_{i=1}^{k} p_i^{a_i}$, the \emph{divisor function} is $\tau(m) = \#\{d \in \mathbb{Z}^+ : d \mid m\}$.
\end{definition}

\begin{theorem}[Divisor Count Formula]
\label{thm:tau}
$\tau(m) = \prod_{i=1}^{k}(a_i + 1)$.
\end{theorem}

\begin{proof}
Each divisor of $m$ has the form $\prod_i p_i^{b_i}$ with $0 \le b_i \le a_i$. The $b_i$ are chosen independently, yielding $\prod_i (a_i+1)$ divisors.
\end{proof}

\begin{theorem}[Multiplicativity]
\label{thm:mult}
If $\gcd(a,b) = 1$, then $\tau(ab) = \tau(a)\,\tau(b)$.
\end{theorem}

\begin{proof}
Coprimality ensures the prime factorizations of $a$ and $b$ involve disjoint prime sets, so the exponents in $ab$ are the union of those in $a$ and $b$. The product formula (Theorem~\ref{thm:tau}) then factors.
\end{proof}

\begin{lemma}[Coprimality of Consecutive Integers]
\label{lem:consec}
For all $n \ge 1$, $\gcd(n, n+1) = 1$.
\end{lemma}

\begin{proof}
If $d \mid n$ and $d \mid (n+1)$, then $d \mid 1$.
\end{proof}

\begin{theorem}[Coprime Factorization of $T_n$]
\label{thm:main}
For all $n \ge 1$:
\[
\tau(T_n) = \begin{cases}
\tau\!\left(\dfrac{n}{2}\right) \cdot \tau(n+1) & \text{if } 2 \mid n, \\[8pt]
\tau(n) \cdot \tau\!\left(\dfrac{n+1}{2}\right) & \text{if } 2 \nmid n.
\end{cases}
\]
\end{theorem}

\begin{proof}
Write $T_n = A \cdot B$ where the factor of $\frac{1}{2}$ is absorbed into the even factor of $n(n+1)$. By Lemma~\ref{lem:consec}, $\gcd(n, n+1)=1$, so every prime dividing $A = n/2$ (when $n$ is even) divides $n$ but not $n+1=B$. Hence $\gcd(A,B) = 1$, and the result follows from Theorem~\ref{thm:mult}.
\end{proof}

\begin{remark}
This decomposition reduces the factoring problem from a single number of magnitude $\sim n^2/2$ to two numbers of magnitude $\sim n$.
\end{remark}

\begin{proposition}[Complexity]
The search algorithm runs in $O(N\sqrt{N})$ time and $O(1)$ auxiliary space, where $N$ is the answer index.
\end{proposition}

\begin{proof}
Each of the $N$ iterations requires trial division of two numbers bounded by $n+1$, costing $O(\sqrt{n})$ each. Summing: $\sum_{n=1}^{N} O(\sqrt{n}) = O(N^{3/2})$.
\end{proof}

\section{Verification}

At $n = 12375$:
\[
T_{12375} = \frac{12375 \times 12376}{2} = 76{,}576{,}500 = 2^2 \cdot 3^2 \cdot 5^3 \cdot 7 \cdot 11 \cdot 13 \cdot 17.
\]
\[
\tau(76{,}576{,}500) = 3 \cdot 3 \cdot 4 \cdot 2 \cdot 2 \cdot 2 \cdot 2 = 576 > 500.
\]

\section{Answer}

\[
\boxed{76576500}
\]

\end{document}
